% TOP_PROBLEM: {"id": 3443, "title": "Outward reloid of composition vs composition of outward reloids", "category_id": 7, "category": {"id": 7, "name": "topology", "display_name": "Topology", "description": "Properties preserved under continuous deformations.", "slug": "topology", "order_index": 7, "created_at": "2026-07-31T15:26:25.671Z"}, "status": "open"}
% FIRST_POSTED: null
% ENTRY_KIND: statement_only
% Imported from: batch04.tex; UnsolvedMath ID 3443
\documentclass[11pt]{article}
\usepackage[margin=25mm]{geometry}
\usepackage{fontspec}
\setmainfont{Latin Modern Roman}
\usepackage{amsmath,amssymb,mathtools}
\usepackage{unicode-math}
\setmathfont{Latin Modern Math}
\usepackage{xurl,hyperref}
\hypersetup{colorlinks=true,urlcolor=blue}
\setcounter{secnumdepth}{0}
\setlength{\parindent}{0pt}
\setlength{\parskip}{5pt}
\setlength{\emergencystretch}{3em}
\newcommand{\sref}[2]{\hyperlink{source:#1:#2}{[S#2]}}
\newcommand{\cataloguescope}[1]{\paragraph{Recorded target.}#1}
\begin{document}
% UnsolvedMath ID 3443; source catalogue code OPG-60020
\setcounter{section}{3442}
\section{Outward reloid of composition vs composition of outward reloids}\label{3443}
{\small\sffamily UnsolvedMath ID 3443 · Topology}
\subsection{Definitions and mathematical statement}

\paragraph{Shared notation.}$\mathbb N=\{1,2,\ldots\}$, and $\mathbb Z,\mathbb Q,\mathbb R,\mathbb C$ denote the integers, rationals, reals and complex numbers. Any other conventions are specified locally.


For a set $X$, let $\mathcal F(X)$ be all filters on $\mathcal P(X)$, including the improper filter $\bot_X=\mathcal P(X)$. A filter is a nonempty family closed under finite intersections and supersets. Order filters by $\mathcal A\sqsubseteq\mathcal B$ when $\mathcal A\supseteq\mathcal B$ as families. Thus $\bigsqcup$ is intersection of filter families and $\bigsqcap$ is the filter generated by their union; the empty join is $\bot_X$ and the empty meet is $\top_X=\{X\}$. For $E\subseteq X$, $\uparrow E=\{U\subseteq X:E\subseteq U\}$ is its principal filter. Write $\mathcal A\mathrel{\#}\mathcal B$ when every $U\in\mathcal A$ and $V\in\mathcal B$ have $U\cap V\ne\varnothing$, equivalently $\mathcal A\sqcap\mathcal B\ne\bot_X$.

A funcoid $f:X\to Y$ is a pair of maps $\alpha_f:\mathcal F(X)\to\mathcal F(Y)$ and $\beta_f:\mathcal F(Y)\to\mathcal F(X)$ satisfying
\[
 \mathcal B\mathrel{\#}\alpha_f(\mathcal A)
 \quad\Longleftrightarrow\quad
 \mathcal A\mathrel{\#}\beta_f(\mathcal B)
 \qquad(\mathcal A\in\mathcal F(X),\ \mathcal B\in\mathcal F(Y)).
\]
Its reverse interchanges $\alpha_f,\beta_f$. Composition satisfies $\alpha_{g\circ f}=\alpha_g\circ\alpha_f$ and $\beta_{g\circ f}=\beta_f\circ\beta_g$. The identity has both maps equal to the identity. The order on $\mathsf{FCD}(X,Y)$ is pointwise order on $\alpha_f$; its meets and joins mean greatest lower bounds and least upper bounds within this complete lattice. A relation $R\subseteq X\times Y$ induces the principal funcoid $\widehat R$ with $\alpha_{\widehat R}(\mathcal A)$ the filter generated by $\{R[U]:U\in\mathcal A\}$ and the analogous inverse-image formula for $\beta$. Here $R[U]=\{y:\exists x\in U,(x,y)\in R\}$.

A reloid $f:X\to Y$ is a filter on $X\times Y$, ordered by reverse inclusion of filter families. For a relation $R\subseteq X\times Y$, its principal reloid is $\uparrow R$. The reloid $g\circ f:X\to Z$ is generated by $G\circ F$ for $F\in f,G\in g$, where $G\circ F=\{(x,z):\exists y,(x,y)\in F,(y,z)\in G\}$. The reverse $f^{-1}$ is generated by $F^{-1}=\{(y,x):(x,y)\in F\}$ for $F\in f$. The identity $1_X$ is $\uparrow\Delta_X$, where $\Delta_X=\{(x,x):x\in X\}$.

For $\mathcal A\in\mathcal F(X),\mathcal B\in\mathcal F(Y)$, $\mathcal A\times^{\mathsf{RLD}}\mathcal B$ is generated by all rectangles $U\times V$ with $U\in\mathcal A,V\in\mathcal B$. Their funcoidal product $\mathcal A\times^{\mathsf{FCD}}\mathcal B$ has first map $\mathcal L\mapsto\mathcal B$ if $\mathcal L\mathrel{\#}\mathcal A$, and $\mathcal L\mapsto\bot_Y$ otherwise; its reverse map has the corresponding formula with $\mathcal A,\mathcal B$ interchanged.

For a funcoid $f:X\to Y$, define its outward reloid and inward reloid by
\[
 O(f)=\bigsqcap^{\mathsf{RLD}}\{\uparrow R:R\subseteq X\times Y,\ f\sqsubseteq\widehat R\},
\]
\[
 I(f)=\bigsqcup^{\mathsf{RLD}}\{\mathcal A\times^{\mathsf{RLD}}\mathcal B:
       \mathcal A\times^{\mathsf{FCD}}\mathcal B\sqsubseteq f\}.
\]
These are the source's $(\mathsf{RLD})_{\mathrm{out}}f$ and $(\mathsf{RLD})_{\mathrm{in}}f$; meets in the first formula generate a filter, and joins in the second intersect filter families.
\paragraph{Statement recorded in the supplied source.}
For all sets $X,Y,Z$ and composable funcoids $f:X\to Y$, $g:Y\to Z$, is
\[
 O(g\circ f)\sqsupseteq O(g)\circ O(f)?
\]
The comparison uses reverse-inclusion order on reloid filter families. No equality assertion or completeness restriction is added. \sref{3443}{2}\sref{3443}{3}
\subsection{Short English statement}
Does the outward reloid of a composite dominate the composite of the two outward reloids?
\subsection{Sources}
\begin{description}
\item[\hypertarget{source:3443:1}{S1}] ULAM.AI and the UnsolvedMath Contributors, UnsolvedMath: A Curated Collection of Open Mathematics Problems, record 3443 (OPG-60020). CC BY 4.0. \url{https://huggingface.co/datasets/ulamai/UnsolvedMath}.
\item[\hypertarget{source:3443:2}{S2}] Open Problem Garden, Outward reloid of composition vs composition of outward reloids, OPG-60020. Original problem statement, full mathematical discussion, bibliography and comments. \url{https://www.openproblemgarden.org/op/outward_reloid_of_composition_vs_composition_of_outward_reloids}.
\item[\hypertarget{source:3443:3}{S3}] Victor Porton, General Topology as Ordered Semigroup Actions, author manuscript, version of 21 July 2026. Reloids induced by a funcoid and composition discussions. \url{https://github.com/vporton/agt-combined/blob/759e45cced7769cf3a84ce60ed6c30f4ce9ffa8e/chap-relationships.tex}.
\item[\hypertarget{source:3443:4}{S4}] Victor Porton, General Topology as Ordered Semigroup Actions, author manuscript, version of 21 July 2026. Basic definitions and composition. \url{https://github.com/vporton/agt-combined/blob/759e45cced7769cf3a84ce60ed6c30f4ce9ffa8e/chap-funcoids.tex}.
\item[\hypertarget{source:3443:5}{S5}] Victor Porton, General Topology as Ordered Semigroup Actions, author manuscript, version of 21 July 2026. Composition of reloids. \url{https://github.com/vporton/agt-combined/blob/759e45cced7769cf3a84ce60ed6c30f4ce9ffa8e/chap-reloids.tex}.
\end{description}
\label{end:3443}
\end{document}
